\chapter{Conclusiones de la etapa}

\textbf{Sobre la implementación.} Más de la mitad del núcleo está hecha, medida en horas, y funciona de extremo a
extremo: registro por API con credencial propia, recepción por MQTT con TLS y por las dos rutas LoRaWAN, pipeline de
seis filtros con cuarentena, supervisión del silencio, alertas por umbral, consulta, exportación y panel en vivo
(capítulo~\ref{cap:pruebas}).

\textbf{Sobre las pruebas.} Los requisitos medibles de esta etapa se cumplen en el banco: 59,7 tramas por segundo sin
pérdida frente a las 50 exigidas, un p95 de 3~ms frente a 2~s, toda caída detectada con 3 alarmas falsas y ninguna
lectura perdida ante caídas del backend y de la base. Además, las pruebas corrigieron el diseño en cinco puntos.

\textbf{Sobre la concordancia y el consumo.} Figuras y modelo de datos se generan desde la documentación y el esquema
real; un verificador de más de 150 comprobaciones falla si informe, documentación o sistema difieren, y cada
comprobación nueva se prueba antes contra el error plantado. Bajo carga, lo construido consume 174~MiB; RNF-17 (2~GB)
se valida en la Etapa~4.

\textbf{Sobre la retroalimentación y lo que sigue.} Las observaciones de la Etapa~1 y del 29 de septiembre están
incorporadas y la entrevista del 5 de octubre adelantó la exportación (sección~\ref{sec:retro}); las de la entrega de
la Etapa~2 se incorporarán al recibirlas. Faltan el sensor físico y la prueba con la pasarela del campus; la Etapa~4
suma la sesión y los roles y valida el sistema en el servidor de referencia.
